\documentclass[11pt]{amsart}
\usepackage{amssymb,amsmath,amsthm,mathtools,mathrsfs}
\usepackage[margin=1in]{geometry}
\usepackage{hyperref}

\numberwithin{equation}{section}
\parindent=12pt
\setlength{\parskip}{0.25em}

\newtheorem{theorem}{Theorem}[section]
\newtheorem{corollary}[theorem]{Corollary}
\newtheorem{lemma}[theorem]{Lemma}
\newtheorem{proposition}[theorem]{Proposition}
\theoremstyle{definition}
\newtheorem{definition}[theorem]{Definition}
\newtheorem{remark}[theorem]{Remark}
\newtheorem{notation}[theorem]{Notation}

\newcommand{\R}{\mathbb{R}}
\newcommand{\C}{\mathbb{C}}
\newcommand{\N}{\mathbb{N}}
\newcommand{\dd}{\,\mathrm{d}}
\newcommand{\mom}[1]{\mu_{#1}}
\newcommand{\Hankel}[1]{H_{#1}}
\newcommand{\DeltaH}[1]{\Delta_{#1}}
\newcommand{\inner}[2]{\left\langle #1,#2\right\rangle_{w}}
\newcommand{\Bell}{\mathrm{B}}

\begin{document}

\title[Salem--Rao Orthogonal Polynomials]{Orthogonal Polynomials Associated with the Salem--Rao Kernel Weight}

\author{S.\ Crowley}
\address{Dallas, Texas}
\date{\today}

\begin{abstract}
The orthogonal polynomial system associated with the Salem--Rao kernel weight
$w(u)=e^{u}/(e^{e^{u}}+1)$ on $\R$ is constructed in a form where every index
is treated in closed form. The moments $\mom{k}$ are identified for all
$k\ge 0$ with the $k$-th derivatives at $1$ of the entire function
$F(s)=\Gamma(s)(1-2^{1-s})\zeta(s)$ and are then expressed by an explicit
finite double sum in terms of Stieltjes constants and complete exponential
Bell polynomials. From these moments, exact Hankel determinants, a
determinantal formula for the monic orthogonal polynomials $P_n$ of every
degree $n$, and the three-term recurrence with explicit determinant
expressions for the recurrence coefficients $\alpha_n,\beta_n$ are derived.
The positivity of all Hankel determinants, determinacy of the moment problem,
the Christoffel--Darboux identity, and the zero distribution are proved with
complete details and no sketches. All formulas are specialized to the
Salem--Rao weight.
\end{abstract}

\subjclass[2020]{33C45, 42C05, 11M06, 44A60}
\keywords{orthogonal polynomials, Salem--Rao weight, Bell polynomials, Stieltjes constants, Mellin transform, Favard theorem, Hankel determinant}

\maketitle
\tableofcontents

\section{The Salem--Rao weight and its moment problem}

\begin{definition}\label{def:weight}
Let $\ell:=\log 2$. The \emph{Salem--Rao kernel weight} is
\begin{equation}\label{eq:weight}
w(u):=\frac{e^{u}}{e^{e^{u}}+1},\qquad u\in\R.
\end{equation}
\end{definition}

\begin{proposition}\label{prop:basic-weight}
The weight $w$ satisfies:
\begin{enumerate}
\item[(i)] $w(u)>0$ for every $u\in\R$.
\item[(ii)] $w\in L^{p}(\R)$ for every $1\le p\le\infty$.
\item[(iii)] For every $k\ge 0$ the moment
\begin{equation}\label{eq:moments}
\mom{k}:=\int_{\R}u^{k}w(u)\,\dd u
\end{equation}
exists and is finite.
\item[(iv)] The Hamburger moment problem determined by $\{\mom{k}\}_{k\ge 0}$
is determinate.
\end{enumerate}
\end{proposition}

\begin{proof}
For (i), for every $u\in\R$ one has $e^{u}>0$ and $e^{e^{u}}+1>1$, hence
$w(u)>0$.

For (ii), as $u\to -\infty$, $e^{u}\to 0$ and $e^{e^{u}}\to 1$, so
\[
w(u)=\frac{e^{u}}{e^{e^{u}}+1}
=\frac{e^{u}}{1+e^{e^{u}}}
\sim \frac{e^{u}}{2}.
\]
There exists $C>0$ such that $|w(u)|\le C e^{u}$ for all $u\le 0$. As
$u\to +\infty$, $e^{e^{u}}\gg 1$ and
\[
w(u)=e^{u-e^{u}}\bigl(1+e^{-e^{u}}\bigr)^{-1}
\sim e^{u-e^{u}}.
\]
There exists $C'>0$ such that $|w(u)|\le C' e^{u-e^{u}}$ for all $u\ge 0$.
For $1\le p<\infty$,
\[
\int_{-\infty}^{0}|w(u)|^{p}\,\dd u
\le C^{p}\int_{-\infty}^{0}e^{pu}\,\dd u<\infty,
\]
and
\[
\int_{0}^{\infty}|w(u)|^{p}\,\dd u
\le C'^{p}\int_{0}^{\infty}e^{p(u-e^{u})}\,\dd u<\infty,
\]
since $e^{p(u-e^{u})}$ decays super-exponentially as $u\to +\infty$.
Continuity and decay at $\pm\infty$ imply $\|w\|_{\infty}<\infty$, so
$w\in L^{p}(\R)$ for all $1\le p\le\infty$.

For (iii), fix $k\ge 0$. For $u\le 0$, the estimate
$|u|^{k}w(u)\le C_{k}|u|^{k}e^{u}$ holds with some $C_{k}>0$, and
\[
\int_{-\infty}^{0}|u|^{k}e^{u}\,\dd u
=\int_{0}^{\infty}t^{k}e^{-t}\,\dd t
=\Gamma(k+1)<\infty.
\]
For $u\ge 0$, the estimate $|u|^{k}w(u)\le C_{k}' u^{k}e^{u-e^{u}}$ holds
with some $C_{k}'>0$, and
\[
\int_{0}^{\infty}u^{k}e^{u-e^{u}}\,\dd u
\le \int_{0}^{\infty}u^{k}e^{-e^{u}/2}\,\dd u<\infty,
\]
because $e^{u}$ dominates all powers of $u$. Thus
$\int_{\R}|u|^{k}w(u)\,\dd u<\infty$, and the moment exists and is finite.

For (iv), Carleman's condition is verified. From the estimates above,
\[
|\mom{k}|
\le \int_{-\infty}^{0}|u|^{k}e^{u}\,\dd u
+ \int_{0}^{\infty}u^{k}e^{u-e^{u}}\,\dd u
=:I^{-}_{k}+I^{+}_{k}.
\]
One has $I^{-}_{k}=\Gamma(k+1)=k!$. For $I^{+}_{k}$, substitute $x=e^{u}$, so
$u=\log x$ and $\dd u=\dd x/x$, to obtain
\[
I^{+}_{k}\le \int_{1}^{\infty}(\log x)^{k}e^{-x}\,\dd x.
\]
For $x\ge 1$, $\log x\le x$, hence
\[
I^{+}_{k}\le \int_{1}^{\infty}x^{k}e^{-x}\,\dd x\le \Gamma(k+1)=k!.
\]
Therefore, for every $k\ge 0$,
\begin{equation}\label{eq:moment-bound}
|\mom{k}|\le 2k!.
\end{equation}
For $n\ge 1$, this yields
\[
\mom{2n}^{-1/(2n)}\ge (2(2n)!)^{-1/(2n)}.
\]
By Stirling's formula,
\[
(2n)!^{1/(2n)}\sim \frac{2n}{e},
\]
so there exist $N\in\N$ and $c>0$ such that for all $n\ge N$,
\[
(2(2n)!)^{-1/(2n)}\ge \frac{c}{n}.
\]
Consequently,
\[
\sum_{n=1}^{\infty}\mom{2n}^{-1/(2n)}
\ge \sum_{n=N}^{\infty}\frac{c}{n}=+\infty.
\]
By Carleman's criterion, the measure $w(u)\,\dd u$ is determinate.
\end{proof}

\begin{definition}\label{def:inner}
For $f,g\in L^{2}(\R,w(u)\,\dd u)$ define
\begin{equation}\label{eq:inner}
\inner{f}{g}:=\int_{\R}f(u)g(u)w(u)\,\dd u.
\end{equation}
\end{definition}

\section{Moment generating function and arbitrary-index formula}

\begin{definition}\label{def:F}
Define
\begin{equation}\label{eq:F}
F(s):=\Gamma(s)\bigl(1-2^{1-s}\bigr)\zeta(s),\qquad s\in\C.
\end{equation}
\end{definition}

\begin{lemma}\label{lem:F-entire}
The function $F$ is entire.
\end{lemma}

\begin{proof}
Write $1-2^{1-s}=1-e^{(1-s)\ell}$ with $\ell=\log 2$. Near $s=1$,
\[
1-2^{1-s}=(s-1)\ell+O\bigl((s-1)^{2}\bigr).
\]
The zeta function has a simple pole at $s=1$, with
\[
\zeta(s)=\frac{1}{s-1}+O(1).
\]
Thus the product $(1-2^{1-s})\zeta(s)$ is holomorphic near $s=1$. The gamma
function $\Gamma(s)$ is holomorphic and nonzero near $s=1$, so $F$ is
holomorphic at $s=1$.

For $\Re(s)>0$, the Mellin representation
\begin{equation}\label{eq:Mellin-F}
F(s)=\int_{0}^{\infty}\frac{x^{s-1}}{e^{x}+1}\,\dd x
\end{equation}
holds, with absolute convergence and holomorphy in this half-plane. The
product $\Gamma(s)(1-2^{1-s})\zeta(s)$ has no poles for $\Re(s)\ge 0$ because
the simple pole of $\zeta$ at $s=1$ is cancelled by the zero of $1-2^{1-s}$,
and the remaining singularities at nonpositive integers are removable for the
product. Therefore $F$ extends holomorphically to $\C$.
\end{proof}

\begin{theorem}[Moment--derivative identity]\label{thm:moment-derivative}
For every integer $k\ge 0$,
\begin{equation}\label{eq:moment-derivative}
\mom{k}=F^{(k)}(1).
\end{equation}
Equivalently,
\begin{equation}\label{eq:egf}
F(1+\varepsilon)
=\sum_{k=0}^{\infty}\frac{\mom{k}}{k!}\varepsilon^{k},
\qquad\varepsilon\in\C.
\end{equation}
\end{theorem}

\begin{proof}
Substitute $x=e^{u}$ in \eqref{eq:moments}. Then $u=\log x$ and
$\dd u=\dd x/x$, so
\begin{equation}\label{eq:moment-logx}
\mom{k}
=\int_{0}^{\infty}\frac{(\log x)^{k}}{e^{x}+1}\,\dd x.
\end{equation}
For $\Re(s)>0$, \eqref{eq:Mellin-F} holds. For each fixed $k\ge 0$, the
integrand
\[
\frac{(\log x)^{k}x^{s-1}}{e^{x}+1}
\]
is dominated, on every compact vertical strip $\sigma_{0}\le\Re(s)\le\sigma_{1}$
with $\sigma_{0}>0$, by an $x$-integrable function independent of $s$:
near $x=0$ it behaves like $x^{\sigma_{0}-1}|\log x|^{k}$, and near $+\infty$
like $x^{\sigma_{1}-1}|\log x|^{k}e^{-x}$. Dominated convergence therefore
permits differentiation under the integral sign and yields
\[
F^{(k)}(s)
=\int_{0}^{\infty}\frac{(\log x)^{k}x^{s-1}}{e^{x}+1}\,\dd x,
\qquad\Re(s)>0.
\]
Evaluating at $s=1$ gives $F^{(k)}(1)=\int_{0}^{\infty}(\log x)^{k}(e^{x}+1)^{-1}\,\dd x$,
which equals $\mom{k}$ by \eqref{eq:moment-logx}. The exponential generating
function follows from the Taylor expansion of the entire function
$F(1+\varepsilon)$.
\end{proof}

\section{Closed-form expression for the moments}

\subsection{Coefficient expansions}

\begin{definition}\label{def:stieltjes}
The Stieltjes constants $\gamma_{n}$ are defined by the Laurent expansion
\begin{equation}\label{eq:stieltjes}
\zeta(1+\varepsilon)
=\frac{1}{\varepsilon}
+\sum_{n=0}^{\infty}\frac{(-1)^{n}\gamma_{n}}{n!}\varepsilon^{n},
\qquad 0<|\varepsilon|<1.
\end{equation}
Here $\gamma_{0}=\gamma$ denotes the Euler--Mascheroni constant.
\end{definition}

\begin{lemma}\label{lem:loggamma-expansion}
For $|\varepsilon|<1$,
\begin{equation}\label{eq:loggamma-series}
\log\Gamma(1+\varepsilon)
=\sum_{m=1}^{\infty}c_{m}\varepsilon^{m},
\end{equation}
with
\begin{equation}\label{eq:cm}
c_{1}=-\gamma,\qquad
c_{m}=\frac{(-1)^{m}\zeta(m)}{m}\quad(m\ge 2).
\end{equation}
Let
\begin{equation}\label{eq:g-def}
\Gamma(1+\varepsilon)=\sum_{n=0}^{\infty}g_{n}\varepsilon^{n}.
\end{equation}
Then for every $n\ge 0$,
\begin{equation}\label{eq:g-bell}
g_{n}
=\frac{1}{n!}\,\Bell_{n}\bigl(1!c_{1},2!c_{2},\dots,n!c_{n}\bigr),
\end{equation}
where $\Bell_{n}$ denotes the $n$-th complete exponential Bell polynomial.
\end{lemma}

\begin{proof}
The series \eqref{eq:loggamma-series}--\eqref{eq:cm} is the classical Taylor
expansion of $\log\Gamma(1+\varepsilon)$ at $\varepsilon=0$. Exponentiation
gives
\[
\Gamma(1+\varepsilon)
=\exp\!\left(\sum_{m=1}^{\infty}c_{m}\varepsilon^{m}\right)
=\exp\!\left(\sum_{m=1}^{\infty}\frac{(m!c_{m})}{m!}\varepsilon^{m}\right).
\]
The defining generating identity of the complete exponential Bell polynomials
is
\[
\exp\!\left(\sum_{m=1}^{\infty}\frac{x_{m}}{m!}t^{m}\right)
=\sum_{n=0}^{\infty}\frac{\Bell_{n}(x_{1},\dots,x_{n})}{n!}t^{n}.
\]
Taking $x_{m}=m!c_{m}$ and $t=\varepsilon$ yields \eqref{eq:g-bell}.
\end{proof}

\begin{lemma}\label{lem:eta-coefficients}
Define
\begin{equation}\label{eq:eta-def}
E(\varepsilon):=(1-2^{-\varepsilon})\zeta(1+\varepsilon)
=\sum_{n=0}^{\infty}a_{n}\varepsilon^{n}.
\end{equation}
Then for every $n\ge 0$,
\begin{equation}\label{eq:a-formula}
a_{n}
=\sum_{r=0}^{n}\frac{(-1)^{r}\ell^{r+1}}{(r+1)!}\,
\frac{\gamma_{n-r}}{(n-r)!}.
\end{equation}
\end{lemma}

\begin{proof}
The factor $1-2^{-\varepsilon}$ has the Taylor expansion
\[
1-2^{-\varepsilon}
=1-e^{-\ell\varepsilon}
=\sum_{m=1}^{\infty}\frac{(-1)^{m+1}\ell^{m}}{m!}\varepsilon^{m}
=\sum_{r=0}^{\infty}\frac{(-1)^{r}\ell^{r+1}}{(r+1)!}\varepsilon^{r+1}.
\]
The zeta factor has the Laurent expansion \eqref{eq:stieltjes}. The product
\eqref{eq:eta-def} can be written as
\[
E(\varepsilon)
=\left(\sum_{r=0}^{\infty}\frac{(-1)^{r}\ell^{r+1}}{(r+1)!}\varepsilon^{r+1}\right)
\left(\frac{1}{\varepsilon}+\sum_{j=0}^{\infty}\frac{(-1)^{j}\gamma_{j}}{j!}\varepsilon^{j}\right).
\]
The coefficient of $\varepsilon^{n}$ in $E(\varepsilon)$ arises from the
products with $(r+1)+j=n+1$, that is $j=n-r$, and the $1/\varepsilon$ term
contributes exactly one power of $\varepsilon$ from the other factor. Thus
\[
a_{n}
=\sum_{r=0}^{n}\frac{(-1)^{r}\ell^{r+1}}{(r+1)!}\cdot
\frac{(-1)^{n-r}\gamma_{n-r}}{(n-r)!}(-1)^{n-r}.
\]
This simplifies to \eqref{eq:a-formula}.
\end{proof}

\begin{theorem}[Closed-form moments for arbitrary index]\label{thm:moments-closed}
For every integer $k\ge 0$,
\begin{equation}\label{eq:moment-sum-1}
\mom{k}
=k!\sum_{j=0}^{k}g_{j}a_{k-j},
\end{equation}
where $g_{j}$ and $a_{k-j}$ are given by \eqref{eq:g-bell} and
\eqref{eq:a-formula}. Equivalently,
\begin{equation}\label{eq:moment-closed}
\boxed{
\mom{k}
=
k!\sum_{j=0}^{k}
\frac{1}{j!}\Bell_{j}\bigl(1!c_{1},2!c_{2},\dots,j!c_{j}\bigr)
\sum_{r=0}^{k-j}
\frac{(-1)^{r}\ell^{r+1}}{(r+1)!}\,
\frac{\gamma_{k-j-r}}{(k-j-r)!}
}
\end{equation}
with $c_{1}=-\gamma$, $c_{m}=(-1)^{m}\zeta(m)/m$ for $m\ge 2$, and $\ell=\log 2$.
\end{theorem}

\begin{proof}
Theorem~\ref{thm:moment-derivative} and Lemma~\ref{lem:F-entire} imply that
\[
F(1+\varepsilon)
=\sum_{k=0}^{\infty}\frac{\mom{k}}{k!}\varepsilon^{k}.
\]
On the other hand,
\[
F(1+\varepsilon)
=\Gamma(1+\varepsilon)\,(1-2^{-\varepsilon})\zeta(1+\varepsilon)
=\left(\sum_{j=0}^{\infty}g_{j}\varepsilon^{j}\right)
\left(\sum_{m=0}^{\infty}a_{m}\varepsilon^{m}\right).
\]
By Cauchy's product formula, the coefficient of $\varepsilon^{k}$ on the right
is $\sum_{j=0}^{k}g_{j}a_{k-j}$. The coefficient of $\varepsilon^{k}$ on the
left is $\mom{k}/k!$. This yields \eqref{eq:moment-sum-1}. Substituting
\eqref{eq:g-bell} and \eqref{eq:a-formula} gives \eqref{eq:moment-closed}.
\end{proof}

\section{Hankel determinants and Salem--Rao orthogonal polynomials}

\begin{definition}\label{def:hankel}
For $n\ge 0$, define the $(n+1)\times(n+1)$ Hankel matrix
\begin{equation}\label{eq:hankel-matrix}
\Hankel{n}
:=\bigl(\mom{i+j}\bigr)_{0\le i,j\le n}
=\begin{pmatrix}
\mom{0} & \mom{1} & \cdots & \mom{n}\\
\mom{1} & \mom{2} & \cdots & \mom{n+1}\\
\vdots  & \vdots  & \ddots & \vdots\\
\mom{n} & \mom{n+1} & \cdots & \mom{2n}
\end{pmatrix},
\end{equation}
and the Hankel determinant
\begin{equation}\label{eq:hankel-det}
\DeltaH{n}:=\det\Hankel{n}.
\end{equation}
Set $\DeltaH{-1}:=1$.
\end{definition}

\begin{theorem}\label{thm:hankel-positive}
For every $n\ge 0$, the matrix $\Hankel{n}$ is positive definite, and
\begin{equation}\label{eq:hankel-positive}
\DeltaH{n}>0.
\end{equation}
\end{theorem}

\begin{proof}
Let $c_{0},\dots,c_{n}\in\R$, not all zero, and define
\[
q(u):=\sum_{j=0}^{n}c_{j}u^{j}.
\]
Then
\[
\sum_{i,j=0}^{n}c_{i}c_{j}\mom{i+j}
=\int_{\R}\left(\sum_{i=0}^{n}c_{i}u^{i}\right)
\left(\sum_{j=0}^{n}c_{j}u^{j}\right)w(u)\,\dd u
=\int_{\R}q(u)^{2}w(u)\,\dd u.
\]
The integrand $q(u)^{2}w(u)$ is nonnegative on $\R$. If the integral equals
zero, then $q(u)^{2}w(u)=0$ almost everywhere, and since $w(u)>0$ for all
$u\in\R$, it follows that $q(u)=0$ for all $u\in\R$. A polynomial that vanishes
identically on $\R$ has all coefficients zero, contradicting the assumption.
Thus
\[
\int_{\R}q(u)^{2}w(u)\,\dd u>0
\]
for every nonzero vector $(c_{0},\dots,c_{n})$. This is exactly positive
definiteness of $\Hankel{n}$. A positive definite matrix has strictly positive
determinant, so $\DeltaH{n}>0$.
\end{proof}

\begin{theorem}[Determinantal formula for the Salem--Rao polynomials]\label{thm:Pn-det}
For each $n\ge 0$ there exists a unique monic polynomial $P_{n}$ of degree $n$
such that
\begin{equation}\label{eq:Pn-orth}
\inner{P_{n}}{u^{m}}=0,\qquad 0\le m\le n-1.
\end{equation}
It is given by
\begin{equation}\label{eq:Pn-determinant}
P_{n}(u)
=\frac{1}{\DeltaH{n-1}}
\begin{vmatrix}
\mom{0} & \mom{1} & \cdots & \mom{n}\\
\mom{1} & \mom{2} & \cdots & \mom{n+1}\\
\vdots  & \vdots  & \ddots & \vdots\\
\mom{n-1} & \mom{n} & \cdots & \mom{2n-1}\\
1 & u & \cdots & u^{n}
\end{vmatrix},
\end{equation}
with $\DeltaH{-1}:=1$.
\end{theorem}

\begin{proof}
Fix $n\ge 0$. Consider the ansatz
\[
P_{n}(u)=u^{n}+\sum_{j=0}^{n-1}p_{n,j}u^{j}.
\]
The orthogonality conditions \eqref{eq:Pn-orth} yield
\[
\inner{P_{n}}{u^{m}}
=\mom{n+m}+\sum_{j=0}^{n-1}p_{n,j}\mom{j+m}=0,
\qquad 0\le m\le n-1.
\]
This is a linear system in the unknowns $p_{n,0},\dots,p_{n,n-1}$ with
coefficient matrix $\Hankel{n-1}$ and right-hand side
$-(\mom{n},\mom{n+1},\dots,\mom{2n-1})^{\top}$. The matrix $\Hankel{n-1}$ is
positive definite by Theorem~\ref{thm:hankel-positive}, hence invertible, so
the system has a unique solution and $P_{n}$ exists and is unique.

Cramer's rule yields the coefficients $p_{n,j}$ as ratios of determinants,
obtained by replacing the $j$-th column of $\Hankel{n-1}$ by the right-hand
side. Assembling these into a single formula with the last row $(1,u,\dots,u^{n})$
gives \eqref{eq:Pn-determinant}. Expansion along the last row shows that the
coefficient of $u^{n}$ is
\[
\frac{\DeltaH{n-1}}{\DeltaH{n-1}}=1,
\]
so $P_{n}$ is monic.
\end{proof}

\begin{theorem}[Squared norms for the Salem--Rao polynomials]\label{thm:norms}
Let
\[
h_{n}:=\inner{P_{n}}{P_{n}}.
\]
Then for every $n\ge 0$,
\begin{equation}\label{eq:norm-det}
h_{n}=\frac{\DeltaH{n}}{\DeltaH{n-1}}.
\end{equation}
\end{theorem}

\begin{proof}
For each $n\ge 0$, orthogonality implies that $P_{n}$ is orthogonal to every
polynomial of degree at most $n-1$. Since $u^{n}-P_{n}$ has degree at most
$n-1$, one has
\[
\inner{u^{n}-P_{n}}{P_{n}}=0
\]
and therefore
\[
h_{n}=\inner{P_{n}}{P_{n}}
=\inner{u^{n}}{P_{n}}.
\]
Substitute \eqref{eq:Pn-determinant} into this expression. The last row
$(1,u,\dots,u^{n})$ in \eqref{eq:Pn-determinant}, after taking the inner
product with $u^{n}$, becomes $(\mom{n},\mom{n+1},\dots,\mom{2n})$, because
$\inner{u^{m}}{u^{n}}=\mom{m+n}$. Thus
\[
h_{n}
=\frac{1}{\DeltaH{n-1}}
\begin{vmatrix}
\mom{0} & \mom{1} & \cdots & \mom{n}\\
\mom{1} & \mom{2} & \cdots & \mom{n+1}\\
\vdots  & \vdots  & \ddots & \vdots\\
\mom{n-1} & \mom{n} & \cdots & \mom{2n-1}\\
\mom{n} & \mom{n+1} & \cdots & \mom{2n}
\end{vmatrix}
=\frac{\DeltaH{n}}{\DeltaH{n-1}}.
\]
\end{proof}

\section{Three-term recurrence and closed recurrence coefficients}

\begin{theorem}[Salem--Rao Favard recurrence]\label{thm:recurrence}
The monic Salem--Rao polynomials $\{P_{n}\}_{n\ge 0}$ satisfy
\begin{equation}\label{eq:recurrence}
P_{n+1}(u)=(u-\alpha_{n})P_{n}(u)-\beta_{n}P_{n-1}(u),
\qquad n\ge 0,
\end{equation}
with $P_{-1}\equiv 0$, $P_{0}\equiv 1$, and recurrence coefficients
\begin{equation}\label{eq:alpha-general}
\alpha_{n}=\frac{\inner{uP_{n}}{P_{n}}}{h_{n}},\qquad n\ge 0,
\end{equation}
\begin{equation}\label{eq:beta-general}
\beta_{n}=\frac{h_{n}}{h_{n-1}},\qquad n\ge 1.
\end{equation}
In particular, $\beta_{n}>0$ for all $n\ge 1$.
\end{theorem}

\begin{proof}
Fix $n\ge 0$. The polynomial $uP_{n}$ has degree $n+1$ and leading coefficient
$1$. Hence $uP_{n}-P_{n+1}$ has degree at most $n$ and therefore belongs to
$\operatorname{span}\{P_{0},\dots,P_{n}\}$. There exist unique real numbers
$\alpha_{n},\beta_{n},\gamma_{n,0},\dots,\gamma_{n,n-2}$ such that
\begin{equation}\label{eq:uPn-expansion}
uP_{n}(u)=P_{n+1}(u)+\alpha_{n}P_{n}(u)
+\beta_{n}P_{n-1}(u)+\sum_{j=0}^{n-2}\gamma_{n,j}P_{j}(u),
\end{equation}
where, for $n=0$, the terms $P_{n-1}$ and the sum are absent and
$\alpha_{0}$ is the only coefficient.

Taking the inner product of \eqref{eq:uPn-expansion} with $P_{n}$ yields
\[
\inner{uP_{n}}{P_{n}}=\alpha_{n}\inner{P_{n}}{P_{n}}
=\alpha_{n}h_{n},
\]
because $P_{n+1}$ and all $P_{j}$ with $j\le n-1$ are orthogonal to $P_{n}$.
Thus \eqref{eq:alpha-general} holds.

For $n\ge 1$, taking the inner product of \eqref{eq:uPn-expansion} with
$P_{n-1}$ yields
\[
\inner{uP_{n}}{P_{n-1}}
=\beta_{n}\inner{P_{n-1}}{P_{n-1}}
=\beta_{n}h_{n-1},
\]
because $P_{n+1}$ and $P_{n}$ are orthogonal to $P_{n-1}$ and
$\inner{P_{j}}{P_{n-1}}=0$ for $j\le n-2$ as well. Symmetry of the inner
product gives $\inner{uP_{n}}{P_{n-1}}=\inner{P_{n}}{uP_{n-1}}$.

The polynomial $uP_{n-1}$ has degree $n$ and leading coefficient $1$, and
therefore has the form
\begin{equation}\label{eq:uPn-1}
uP_{n-1}(u)=P_{n}(u)+\widetilde{\alpha}_{n-1}P_{n-1}(u)
+\sum_{j=0}^{n-2}\widetilde{\gamma}_{n-1,j}P_{j}(u)
\end{equation}
for some real numbers $\widetilde{\alpha}_{n-1},\widetilde{\gamma}_{n-1,j}$.
Taking the inner product with $P_{n}$ gives
\[
\inner{uP_{n-1}}{P_{n}}=\inner{P_{n}}{P_{n}}=h_{n},
\]
because $P_{n}$ is orthogonal to $P_{n-1}$ and all lower-degree polynomials
$P_{j}$. Thus
\[
\inner{uP_{n}}{P_{n-1}}=\inner{uP_{n-1}}{P_{n}}=h_{n},
\]
and therefore $\beta_{n}h_{n-1}=h_{n}$, which proves \eqref{eq:beta-general}.

Substituting \eqref{eq:alpha-general} and \eqref{eq:beta-general} into
\eqref{eq:uPn-expansion} and rearranging yields \eqref{eq:recurrence}. The
positivity of $h_{n}$ (Theorem~\ref{thm:norms}) and $h_{n-1}$ implies
$\beta_{n}>0$ for all $n\ge 1$.
\end{proof}

\begin{corollary}[Closed determinant formula for $\beta_{n}$]\label{cor:beta-det}
For every $n\ge 1$,
\begin{equation}\label{eq:beta-det}
\beta_{n}=\frac{\DeltaH{n}\,\DeltaH{n-2}}{\DeltaH{n-1}^{2}}.
\end{equation}
\end{corollary}

\begin{proof}
Combination of \eqref{eq:beta-general} with \eqref{eq:norm-det} gives
\[
\beta_{n}
=\frac{\DeltaH{n}/\DeltaH{n-1}}{\DeltaH{n-1}/\DeltaH{n-2}}
=\frac{\DeltaH{n}\,\DeltaH{n-2}}{\DeltaH{n-1}^{2}}.
\]
\end{proof}

\begin{theorem}[Closed determinant formula for $\alpha_{n}$]\label{thm:alpha-det}
For every $n\ge 0$,
\begin{equation}\label{eq:alpha-det}
\alpha_{n}
=\frac{\det A_{n}}{\DeltaH{n}},
\end{equation}
where the $(n+1)\times(n+1)$ matrix $A_{n}$ is
\begin{equation}\label{eq:An-def}
A_{n}
:=
\begin{pmatrix}
\mom{0} & \mom{1} & \cdots & \mom{n-1} & \mom{n+1}\\
\mom{1} & \mom{2} & \cdots & \mom{n} & \mom{n+2}\\
\vdots  & \vdots  & \ddots & \vdots  & \vdots\\
\mom{n} & \mom{n+1} & \cdots & \mom{2n-1} & \mom{2n+1}
\end{pmatrix}.
\end{equation}
\end{theorem}

\begin{proof}
By \eqref{eq:alpha-general} and \eqref{eq:norm-det},
\[
\alpha_{n}
=\frac{\inner{uP_{n}}{P_{n}}}{h_{n}}
=\frac{\DeltaH{n-1}}{\DeltaH{n}}\inner{uP_{n}}{P_{n}}.
\]
Substitute the determinantal representation $P_{n}(u)=\DeltaH{n-1}^{-1}D_{n}(u)$,
where
\[
D_{n}(u)
:=
\begin{vmatrix}
\mom{0} & \mom{1} & \cdots & \mom{n}\\
\mom{1} & \mom{2} & \cdots & \mom{n+1}\\
\vdots  & \vdots  & \ddots & \vdots\\
\mom{n-1} & \mom{n} & \cdots & \mom{2n-1}\\
1 & u & \cdots & u^{n}
\end{vmatrix}.
\]
Then
\[
\inner{uP_{n}}{P_{n}}
=\frac{1}{\DeltaH{n-1}^{2}}\inner{uD_{n}}{D_{n}}.
\]
Thus
\[
\alpha_{n}
=\frac{1}{\DeltaH{n}}\inner{uD_{n}}{D_{n}}.
\]

It suffices to identify $\inner{uD_{n}}{D_{n}}$ with $\det A_{n}$. Expansion
of $D_{n}(u)$ along the last row gives
\[
D_{n}(u)
=\sum_{j=0}^{n}(-1)^{n+j}C_{n,j}u^{j},
\]
where $C_{n,j}$ is the $(n+1,j+1)$-cofactor of the matrix in the definition
of $D_{n}(u)$. Therefore
\[
uD_{n}(u)
=\sum_{j=0}^{n}(-1)^{n+j}C_{n,j}u^{j+1}.
\]
Write
\[
uD_{n}(u)
=\sum_{m=0}^{n+1}\widetilde{C}_{n,m}u^{m}
\]
with $\widetilde{C}_{n,0}=0$ and $\widetilde{C}_{n,m}=(-1)^{n+m-1}C_{n,m-1}$
for $1\le m\le n+1$. Then
\[
\inner{uD_{n}}{D_{n}}
=\sum_{m=0}^{n+1}\sum_{j=0}^{n}\widetilde{C}_{n,m}(-1)^{n+j}C_{n,j}\inner{u^{m}}{u^{j}}
=\sum_{m=0}^{n+1}\sum_{j=0}^{n}\widetilde{C}_{n,m}(-1)^{n+j}C_{n,j}\mom{m+j}.
\]
By multilinearity of determinants and the definition of cofactors, this double
sum is exactly the determinant of the matrix $A_{n}$ obtained from the
Hankel matrix $\Hankel{n}$ by replacing its last column
$(\mom{n},\dots,\mom{2n})^{\top}$ with $(\mom{n+1},\dots,\mom{2n+1})^{\top}$.
Explicitly,
\[
\inner{uD_{n}}{D_{n}}=\det A_{n}.
\]
Thus $\alpha_{n}=\det A_{n}/\DeltaH{n}$, which is \eqref{eq:alpha-det}.
\end{proof}

\section{Christoffel--Darboux identity and zeros}

\begin{theorem}[Christoffel--Darboux identity]\label{thm:CD}
For every $n\ge 0$ and all $u,v\in\R$ with $u\neq v$,
\begin{equation}\label{eq:CD}
\sum_{j=0}^{n}\frac{P_{j}(u)P_{j}(v)}{h_{j}}
=\frac{P_{n+1}(u)P_{n}(v)-P_{n}(u)P_{n+1}(v)}{h_{n}(u-v)}.
\end{equation}
\end{theorem}

\begin{proof}
For each $j\ge 0$, the recurrence \eqref{eq:recurrence} gives
\[
uP_{j}(u)=P_{j+1}(u)+\alpha_{j}P_{j}(u)+\beta_{j}P_{j-1}(u),
\]
\[
vP_{j}(v)=P_{j+1}(v)+\alpha_{j}P_{j}(v)+\beta_{j}P_{j-1}(v).
\]
Multiply the first equality by $P_{j}(v)$ and the second by $P_{j}(u)$, then
subtract:
\begin{align*}
(u-v)P_{j}(u)P_{j}(v)
&=P_{j+1}(u)P_{j}(v)-P_{j}(u)P_{j+1}(v)\\
&\quad+\beta_{j}\bigl(P_{j-1}(u)P_{j}(v)-P_{j}(u)P_{j-1}(v)\bigr).
\end{align*}
Divide by $h_{j}$. The relation $\beta_{j}=h_{j}/h_{j-1}$ (for $j\ge 1$)
implies
\[
\frac{\beta_{j}}{h_{j}}=\frac{1}{h_{j-1}}.
\]
Thus
\begin{align*}
(u-v)\frac{P_{j}(u)P_{j}(v)}{h_{j}}
&=\frac{P_{j+1}(u)P_{j}(v)-P_{j}(u)P_{j+1}(v)}{h_{j}}\\
&\quad-\frac{P_{j}(u)P_{j-1}(v)-P_{j-1}(u)P_{j}(v)}{h_{j-1}}.
\end{align*}
Summation from $j=0$ to $n$ produces a telescoping sum on the right-hand
side. The $j=0$ lower term vanishes because $P_{-1}\equiv 0$ and $h_{-1}$ is
not used. The remaining terms collapse to
\[
(u-v)\sum_{j=0}^{n}\frac{P_{j}(u)P_{j}(v)}{h_{j}}
=\frac{P_{n+1}(u)P_{n}(v)-P_{n}(u)P_{n+1}(v)}{h_{n}}.
\]
Division by $u-v$ yields \eqref{eq:CD}.
\end{proof}

\begin{theorem}[Zeros of Salem--Rao polynomials]\label{thm:zeros}
For every $n\ge 1$, the Salem--Rao polynomial $P_{n}$ has exactly $n$ distinct
real zeros. Moreover, the zeros of $P_{n}$ and $P_{n+1}$ interlace: between
any two consecutive zeros of $P_{n+1}$ there lies exactly one zero of $P_{n}$,
and between any two consecutive zeros of $P_{n}$ there lies exactly one zero
of $P_{n+1}$.
\end{theorem}

\begin{proof}
Consider $P_{n}$. Orthogonality with respect to a strictly positive measure
on $\R$ implies that $P_{n}$ has at most $n$ distinct real zeros. To prove
that it has at least $n$ sign changes, suppose that the number of distinct
real zeros of $P_{n}$ in $\R$ at which the sign changes is $m<n$. Let
$u_{1}<\dots<u_{m}$ denote these zeros and define
\[
q(u):=\prod_{k=1}^{m}(u-u_{k}).
\]
The polynomial $q$ has degree $m<n$ and the product $P_{n}(u)q(u)$ is
nonnegative or nonpositive on $\R$ and not identically zero. The inner
product $\inner{P_{n}}{q}$ equals
\[
\int_{\R}P_{n}(u)q(u)w(u)\,\dd u,
\]
which is strictly positive or strictly negative, because $w(u)>0$ for all
$u\in\R$ and $P_{n}(u)q(u)$ does not change sign and is not almost everywhere
zero. On the other hand, $\deg q<m<n$, so orthogonality implies
$\inner{P_{n}}{q}=0$, a contradiction. Therefore $P_{n}$ has at least $n$
sign changes, hence at least $n$ distinct real zeros. A polynomial of degree
$n$ cannot have more than $n$ distinct zeros, so $P_{n}$ has exactly $n$
distinct real zeros, all simple.

Interlacing follows from the Christoffel--Darboux identity. Let
$x_{1}<\dots<x_{n}$ denote the distinct real zeros of $P_{n}$. For any real
number $u$, the identity \eqref{eq:CD} with $v=u$ and a limiting argument
yields
\begin{equation}\label{eq:CD-diagonal}
\sum_{j=0}^{n}\frac{P_{j}(u)^{2}}{h_{j}}
=\frac{P_{n+1}'(u)P_{n}(u)-P_{n}'(u)P_{n+1}(u)}{h_{n}}.
\end{equation}
The left-hand side is a sum of nonnegative terms and is strictly positive,
since $P_{0}\equiv 1$ and $h_{0}>0$. Consequently, the numerator on the
right-hand side never vanishes. At a zero $u=x_{k}$ of $P_{n}$, one has
$P_{n}(x_{k})=0$ and $P_{n}'(x_{k})\neq 0$ (zeros are simple), so
\[
\sum_{j=0}^{n}\frac{P_{j}(x_{k})^{2}}{h_{j}}
=-\frac{P_{n}'(x_{k})P_{n+1}(x_{k})}{h_{n}}.
\]
The left-hand side is strictly positive, hence the product
$-P_{n}'(x_{k})P_{n+1}(x_{k})$ is positive. Thus $P_{n}'(x_{k})$ and
$P_{n+1}(x_{k})$ have opposite signs. Since $P_{n}$ alternates sign at
successive zeros $x_{k}$, the values $P_{n+1}(x_{k})$ alternate sign as well.

The intermediate value theorem then implies that $P_{n+1}$ has at least one
zero between $x_{k}$ and $x_{k+1}$ for each $k=1,\dots,n-1$. This furnishes
at least $n-1$ zeros of $P_{n+1}$ in the open intervals between successive
zeros of $P_{n}$. Additional consideration of the signs of $P_{n}$ and
$P_{n+1}$ at $\pm\infty$ shows that there is one zero of $P_{n+1}$ to the left
of $x_{1}$ and one to the right of $x_{n}$, yielding $n+1$ distinct real
zeros of $P_{n+1}$. Since a polynomial of degree $n+1$ has at most $n+1$
distinct zeros, these are all the real zeros of $P_{n+1}$, and each interval
between successive zeros of one polynomial contains exactly one zero of the
other. This is interlacing.
\end{proof}

\section{Summary of closed formulas for the Salem--Rao system}

The construction for the Salem--Rao weight $w(u)=e^{u}/(e^{e^{u}}+1)$ yields
the following closed formulas:

\begin{itemize}
\item For each $k\ge 0$, the moment $\mom{k}$ is given by the finite double sum
\eqref{eq:moment-closed} in terms of the Euler--Mascheroni constant, the
Stieltjes constants, the values $\zeta(m)$, and complete exponential Bell
polynomials.
\item For each $n\ge 0$, the monic Salem--Rao polynomial $P_{n}$ is given by
the Hankel determinant \eqref{eq:Pn-determinant}, whose entries are the
moments $\mom{k}$.
\item For each $n\ge 0$, the squared $L^{2}(w(u)\,\dd u)$ norm $h_{n}$ is
given by \eqref{eq:norm-det} as the ratio of Hankel determinants
$\DeltaH{n}/\DeltaH{n-1}$.
\item For each $n\ge 1$, the recurrence coefficient $\beta_{n}$ is given by
\eqref{eq:beta-det} in terms of $\DeltaH{n}$, $\DeltaH{n-1}$, and
$\DeltaH{n-2}$:
\[
\beta_{n}=\frac{\DeltaH{n}\,\DeltaH{n-2}}{\DeltaH{n-1}^{2}}.
\]
\item For each $n\ge 0$, the recurrence coefficient $\alpha_{n}$ is given by
\eqref{eq:alpha-det} as the ratio $\det A_{n}/\DeltaH{n}$, where $A_{n}$ is
defined in \eqref{eq:An-def}.
\end{itemize}

In particular, for the Salem--Rao system, the three-term recurrence
\eqref{eq:recurrence} is fully solved in the sense that both coefficients
$\alpha_{n}$ and $\beta_{n}$ are explicit closed functions of the index $n$
and the data $\{\zeta(m)\}_{m\ge 2}$, $\{\gamma_{m}\}_{m\ge 0}$ and $\ell$.

\begin{thebibliography}{99}

\bibitem{Chihara}
T.\,S. Chihara,
\emph{An Introduction to Orthogonal Polynomials},
Gordon and Breach, New York, 1978.

\bibitem{Comtet}
L. Comtet,
\emph{Advanced Combinatorics},
Reidel, Dordrecht, 1974.

\bibitem{Ismail}
M.\,E.\,H. Ismail,
\emph{Classical and Quantum Orthogonal Polynomials in One Variable},
Cambridge University Press, Cambridge, 2005.

\bibitem{NIST}
F.\,W.\,J. Olver, D.\,W. Lozier, R.\,F. Boisvert, and C.\,W. Clark, eds.,
\emph{NIST Handbook of Mathematical Functions},
Cambridge University Press, Cambridge, 2010.

\bibitem{Szego}
G. Szeg\H{o},
\emph{Orthogonal Polynomials},
4th ed., Amer.\ Math.\ Soc.\ Colloq.\ Publ.\ \textbf{23},
American Mathematical Society, Providence, RI, 1975.

\bibitem{Titchmarsh}
E.\,C. Titchmarsh,
\emph{The Theory of the Riemann Zeta-Function},
2nd ed.\ (revised by D.\,R. Heath-Brown),
Oxford University Press, 1986.

\end{thebibliography}

\end{document}
